% Options for packages loaded elsewhere
% Options for packages loaded elsewhere
\PassOptionsToPackage{unicode}{hyperref}
\PassOptionsToPackage{hyphens}{url}
\PassOptionsToPackage{dvipsnames,svgnames,x11names}{xcolor}
%
\documentclass[
  12pt,
]{article}
\usepackage{xcolor}
\usepackage[margin=1in]{geometry}
\usepackage{amsmath,amssymb}
\setcounter{secnumdepth}{-\maxdimen} % remove section numbering
\usepackage{iftex}
\ifPDFTeX
  \usepackage[T1]{fontenc}
  \usepackage[utf8]{inputenc}
  \usepackage{textcomp} % provide euro and other symbols
\else % if luatex or xetex
  \usepackage{unicode-math} % this also loads fontspec
  \defaultfontfeatures{Scale=MatchLowercase}
  \defaultfontfeatures[\rmfamily]{Ligatures=TeX,Scale=1}
\fi
\usepackage{lmodern}
\ifPDFTeX\else
  % xetex/luatex font selection
\fi
% Use upquote if available, for straight quotes in verbatim environments
\IfFileExists{upquote.sty}{\usepackage{upquote}}{}
\IfFileExists{microtype.sty}{% use microtype if available
  \usepackage[]{microtype}
  \UseMicrotypeSet[protrusion]{basicmath} % disable protrusion for tt fonts
}{}
\usepackage{setspace}
\makeatletter
\@ifundefined{KOMAClassName}{% if non-KOMA class
  \IfFileExists{parskip.sty}{%
    \usepackage{parskip}
  }{% else
    \setlength{\parindent}{0pt}
    \setlength{\parskip}{6pt plus 2pt minus 1pt}}
}{% if KOMA class
  \KOMAoptions{parskip=half}}
\makeatother
% Make \paragraph and \subparagraph free-standing
\makeatletter
\ifx\paragraph\undefined\else
  \let\oldparagraph\paragraph
  \renewcommand{\paragraph}{
    \@ifstar
      \xxxParagraphStar
      \xxxParagraphNoStar
  }
  \newcommand{\xxxParagraphStar}[1]{\oldparagraph*{#1}\mbox{}}
  \newcommand{\xxxParagraphNoStar}[1]{\oldparagraph{#1}\mbox{}}
\fi
\ifx\subparagraph\undefined\else
  \let\oldsubparagraph\subparagraph
  \renewcommand{\subparagraph}{
    \@ifstar
      \xxxSubParagraphStar
      \xxxSubParagraphNoStar
  }
  \newcommand{\xxxSubParagraphStar}[1]{\oldsubparagraph*{#1}\mbox{}}
  \newcommand{\xxxSubParagraphNoStar}[1]{\oldsubparagraph{#1}\mbox{}}
\fi
\makeatother


\usepackage{longtable,booktabs,array}
\newcounter{none} % for unnumbered tables
\usepackage{calc} % for calculating minipage widths
% Correct order of tables after \paragraph or \subparagraph
\usepackage{etoolbox}
\makeatletter
\patchcmd\longtable{\par}{\if@noskipsec\mbox{}\fi\par}{}{}
\makeatother
% Allow footnotes in longtable head/foot
\IfFileExists{footnotehyper.sty}{\usepackage{footnotehyper}}{\usepackage{footnote}}
\makesavenoteenv{longtable}
\usepackage{graphicx}
\makeatletter
\newsavebox\pandoc@box
\newcommand*\pandocbounded[1]{% scales image to fit in text height/width
  \sbox\pandoc@box{#1}%
  \Gscale@div\@tempa{\textheight}{\dimexpr\ht\pandoc@box+\dp\pandoc@box\relax}%
  \Gscale@div\@tempb{\linewidth}{\wd\pandoc@box}%
  \ifdim\@tempb\p@<\@tempa\p@\let\@tempa\@tempb\fi% select the smaller of both
  \ifdim\@tempa\p@<\p@\scalebox{\@tempa}{\usebox\pandoc@box}%
  \else\usebox{\pandoc@box}%
  \fi%
}
% Set default figure placement to htbp
\def\fps@figure{htbp}
\makeatother


% definitions for citeproc citations
\NewDocumentCommand\citeproctext{}{}
\NewDocumentCommand\citeproc{mm}{%
  \begingroup\def\citeproctext{#2}\cite{#1}\endgroup}
\makeatletter
 % allow citations to break across lines
 \let\@cite@ofmt\@firstofone
 % avoid brackets around text for \cite:
 \def\@biblabel#1{}
 \def\@cite#1#2{{#1\if@tempswa , #2\fi}}
\makeatother
\newlength{\cslhangindent}
\setlength{\cslhangindent}{1.5em}
\newlength{\csllabelwidth}
\setlength{\csllabelwidth}{3em}
\newenvironment{CSLReferences}[2] % #1 hanging-indent, #2 entry-spacing
 {\begin{list}{}{%
  \setlength{\itemindent}{0pt}
  \setlength{\leftmargin}{0pt}
  \setlength{\parsep}{0pt}
  % turn on hanging indent if param 1 is 1
  \ifodd #1
   \setlength{\leftmargin}{\cslhangindent}
   \setlength{\itemindent}{-1\cslhangindent}
  \fi
  % set entry spacing
  \setlength{\itemsep}{#2\baselineskip}}}
 {\end{list}}
\usepackage{calc}
\newcommand{\CSLBlock}[1]{\hfill\break\parbox[t]{\linewidth}{\strut\ignorespaces#1\strut}}
\newcommand{\CSLLeftMargin}[1]{\parbox[t]{\csllabelwidth}{\strut#1\strut}}
\newcommand{\CSLRightInline}[1]{\parbox[t]{\linewidth - \csllabelwidth}{\strut#1\strut}}
\newcommand{\CSLIndent}[1]{\hspace{\cslhangindent}#1}



\setlength{\emergencystretch}{3em} % prevent overfull lines

\providecommand{\tightlist}{%
  \setlength{\itemsep}{0pt}\setlength{\parskip}{0pt}}



 


\makeatletter
\@ifpackageloaded{caption}{}{\usepackage{caption}}
\AtBeginDocument{%
\ifdefined\contentsname
  \renewcommand*\contentsname{Table of contents}
\else
  \newcommand\contentsname{Table of contents}
\fi
\ifdefined\listfigurename
  \renewcommand*\listfigurename{List of Figures}
\else
  \newcommand\listfigurename{List of Figures}
\fi
\ifdefined\listtablename
  \renewcommand*\listtablename{List of Tables}
\else
  \newcommand\listtablename{List of Tables}
\fi
\ifdefined\figurename
  \renewcommand*\figurename{Figure}
\else
  \newcommand\figurename{Figure}
\fi
\ifdefined\tablename
  \renewcommand*\tablename{Table}
\else
  \newcommand\tablename{Table}
\fi
}
\@ifpackageloaded{float}{}{\usepackage{float}}
\floatstyle{ruled}
\@ifundefined{c@chapter}{\newfloat{codelisting}{h}{lop}}{\newfloat{codelisting}{h}{lop}[chapter]}
\floatname{codelisting}{Listing}
\newcommand*\listoflistings{\listof{codelisting}{List of Listings}}
\makeatother
\makeatletter
\makeatother
\makeatletter
\@ifpackageloaded{caption}{}{\usepackage{caption}}
\@ifpackageloaded{subcaption}{}{\usepackage{subcaption}}
\makeatother
\usepackage{bookmark}
\IfFileExists{xurl.sty}{\usepackage{xurl}}{} % add URL line breaks if available
\urlstyle{same}
\hypersetup{
  pdftitle={Decoding Pancreatic Ductal Adenocarcinoma Drivers via Consensus Meta-Analysis: A Prognostic Gene Signature Resilient to Stromal Contamination},
  pdfauthor={Md. Jubayer Hossain; Md. Shakil Ahmed; Muntasim Fuad; Muhibullah Shahjahan},
  pdfkeywords={pancreatic cancer, PDAC, meta-analysis, desmoplastic
stroma, OASL, EMT, prognostic signature, LASSO-Cox},
  colorlinks=true,
  linkcolor={blue},
  filecolor={Maroon},
  citecolor={Blue},
  urlcolor={Blue},
  pdfcreator={LaTeX via pandoc}}


\title{Decoding Pancreatic Ductal Adenocarcinoma Drivers via Consensus
Meta-Analysis: A Prognostic Gene Signature Resilient to Stromal
Contamination}
\author{Md. Jubayer Hossain \and Md. Shakil Ahmed \and Muntasim
Fuad \and Muhibullah Shahjahan}
\date{2026-05-20}
\begin{document}
\maketitle
\begin{abstract}
Pancreatic ductal adenocarcinoma (PDAC) is among the most lethal solid
tumors, with transcriptomic characterization complicated by the abundant
desmoplastic stroma that constitutes up to 90\% of tumor volume. We
performed a random-effects meta-analysis of seven independent GEO
RNA-sequencing cohorts to identify a consensus set of study-invariant
PDAC drivers. The stringent meta-analytic framework identified 418
consensus differentially expressed genes (DEGs; rem\_padj \textless{}
0.05, \textbar pooled log₂FC\textbar{} ≥ 1)---a conservative signature
reflecting genes that consistently emerge despite cross-cohort
variability in tumor purity and stromal composition. Functional
enrichment revealed upregulation of epithelial-mesenchymal transition
(EMT; NES = +2.48), interferon-gamma response (NES = +2.13), and
allograft rejection (NES = +1.98) pathways, implicating stromal
activation and immune exclusion as defining features of the consensus
PDAC transcriptome. Hub gene analysis identified FN1, COL1A1, COL1A2,
COL3A1, and MMP9 as the most topologically central nodes---all
components of the desmoplastic stroma. TCGA PAAD validation yielded a
non-significant Spearman correlation (r = +0.064, p = 0.31), consistent
with known variable tumor purity in TCGA PDAC specimens relative to the
GEO cohorts. Nineteen genes achieved FDR-significant prognostic
associations in TCGA PAAD (n = 177, 93 events), with OASL (HR = 1.346,
padj = 0.001), OAS3 (HR = 1.506), and TGFBI (HR = 1.359) as the
strongest individual predictors. A complementary 11-gene LASSO-Cox risk
model achieved a concordance index of 0.666 (log-rank p = 9.35 × 10⁻⁵).
These findings establish a stringently filtered transcriptomic resource
for PDAC that captures the interplay between tumor-intrinsic programs
and the stromal microenvironment.
\end{abstract}


\setstretch{2}
\section{Introduction}\label{introduction}

Pancreatic ductal adenocarcinoma (PDAC) remains one of the most lethal
malignancies, with a five-year survival rate of approximately 12\% and a
median overall survival of 6--12 months in the metastatic setting
(Siegel et al. 2023). Unlike many solid tumors where advances in
targeted therapy and immunotherapy have substantially improved outcomes,
PDAC has proven refractory to most systemic treatments, and
gemcitabine-based chemotherapy remains a mainstay of palliative
management (Conroy et al. 2011). The discovery of robust prognostic
biomarkers and the identification of novel therapeutic targets represent
pressing unmet needs.

A defining pathological feature of PDAC is its dense desmoplastic
stroma, which may constitute 50--90\% of total tumor volume and consists
of activated pancreatic stellate cells (PSCs), cancer-associated
fibroblasts (CAFs), extracellular matrix components, and various immune
cell populations (Feig et al. 2013). This stromal abundance poses a
fundamental challenge for transcriptomic biomarker discovery: bulk
RNA-seq measurements reflect a mixture of tumor-cell-intrinsic and
stromal signals, and the relative contribution of each varies across
specimens depending on tumor purity, biopsy site, and processing
methodology. Individual-cohort studies of PDAC transcriptomics are
therefore particularly susceptible to reproducibility failures arising
from cohort-specific differences in tumor-stroma composition (Moffitt et
al. 2015).

Meta-analytic integration across multiple independent PDAC cohorts
offers a principled approach to this challenge. Genes whose differential
expression is consistently detected across multiple cohorts with varying
tumor purity---a meta-analysis of consensus DEGs---are more likely to
represent genuine tumor-intrinsic or consistently co-regulated stromal
signals than genes identified in only one or two studies. The smaller
number of DEGs expected from a stringent multi-cohort PDAC meta-analysis
(compared to cancers with less stromal contamination) is itself
informative: it reflects the higher biological standard required for
cross-cohort reproducibility in a heterogeneous tumor type.

The present study performs a seven-cohort random-effects meta-analysis
of PDAC transcriptomics, yielding 418 consensus DEGs that represent the
study-invariant molecular signature of PDAC. We characterize this
signature through functional enrichment, PPI network analysis, TCGA
validation, survival analysis, and LASSO-Cox risk score modeling,
interpreting each finding in the context of PDAC's defining
biology---tumor-stroma crosstalk and immune microenvironment remodeling.

\section{Materials and Methods}\label{materials-and-methods}

\subsection{Dataset Acquisition and
Preprocessing}\label{dataset-acquisition-and-preprocessing}

Seven GEO RNA-seq datasets were analyzed: GSE130688, GSE136569,
GSE171485, GSE196009, GSE211398, GSE280271, and GSE293744. All datasets
included paired pancreatic ductal adenocarcinoma and adjacent normal
pancreatic tissue. Preprocessing used DESeq2 VST normalization (Love et
al. 2014). REML random-effects meta-analysis was performed per gene
using metafor (Viechtbauer 2010). Consensus DEGs: rem\_padj \textless{}
0.05 and \textbar pooled log₂FC\textbar{} ≥ 1.

\subsection{TCGA Validation}\label{tcga-validation}

TCGA PAAD data were retrieved via curatedTCGAData (v2.1.1,
RNASeq2GeneNorm; (Ramos et al. 2017)). Spearman correlation was computed
between meta-analysis pooled log₂FC and TCGA log₂FC. Given the variable
tumor purity in TCGA PAAD specimens, we also examined concordance at the
gene-set level (hub gene enrichment).

\subsection{Functional Enrichment, Network, Survival, and
ML}\label{functional-enrichment-network-survival-and-ml}

Methods followed the pipeline described in the kidney Methods. GSEA
Hallmark, ORA KEGG/GO, STRINGdb (score ≥ 700) (Szklarczyk et al. 2023),
and clusterProfiler (Wu et al. 2021) were employed. Survival analysis
used TCGA PAAD OS data (n = 177; 93 events). LASSO-Cox (α = 0.9, 10-fold
CV; (Friedman et al. 2010)) was fitted using hub genes as features.
Random Forest was not performed due to insufficient TCGA PAAD normal
tissue samples (n \textless{} 5), which precluded reliable stratified
cross-validation; this limitation is explicitly acknowledged.

\section{Results}\label{results}

\subsection{Meta-Analysis of Seven Cohorts Yields 418 Study-Invariant
Consensus
DEGs}\label{meta-analysis-of-seven-cohorts-yields-418-study-invariant-consensus-degs}

Integration of seven GEO PDAC RNA-seq datasets through REML
random-effects meta-analysis identified 418 consensus DEGs (rem\_padj
\textless{} 0.05; \textbar pooled log₂FC\textbar{} ≥ 1). The relatively
small number of consensus DEGs compared to other cancer types analyzed
in this study series (kidney: 2,766; colorectal: 2,170; esophagus:
2,082) reflects the biological and technical challenges of PDAC
transcriptomics: the abundant desmoplastic stroma introduces
cohort-specific variability in measured transcript levels, and the more
stringent inter-cohort consistency required by the random-effects
framework necessarily reduces the number of surviving DEGs in a tumor
type with high compositional heterogeneity. Conversely, the 418 genes
that do pass this stringent threshold represent a high-confidence set of
study-invariant PDAC molecular features, more likely to reflect genuine
biological signals than the larger gene sets identified in less
heterogeneous cancers.

All 418 consensus DEGs were upregulated (positive pooled log₂FC), with
the top-ranking genes including ECM components and immune-response
effectors. The absence of significant downregulated DEGs at this
threshold may reflect the dilution of tumor-intrinsic suppression events
by stromal signal, which predominates in bulk tissue measurements.

\subsection{Functional Enrichment Reveals EMT and Immune Activation
Signatures}\label{functional-enrichment-reveals-emt-and-immune-activation-signatures}

GSEA Hallmark analysis identified three gene sets with the highest NES
values, all representing upregulation in PDAC tumors relative to normal
pancreas (Figure 3). HALLMARK\_EPITHELIAL\_MESENCHYMAL\_TRANSITION (NES
= +2.48, FDR \textless{} 10⁻⁹) ranked first, reflecting the prominent
mesenchymal reprogramming that characterizes PDAC tumor cells and the
activated stroma of cancer-associated fibroblasts and PSCs (Kalluri
2016). EMT is a central mechanism driving PDAC invasiveness and
treatment resistance, and its identification as the dominant GSEA
enrichment signal confirms that the meta-analysis consensus signature
captures the most biologically meaningful aspects of PDAC pathobiology.
HALLMARK\_INTERFERON\_GAMMA\_RESPONSE (NES = +2.13) and
HALLMARK\_ALLOGRAFT\_REJECTION (NES = +1.98) ranked second and third,
respectively, reflecting the chronic inflammatory state that
characterizes the PDAC microenvironment, where persistent innate and
adaptive immune activation creates a paradoxically immunosuppressive
milieu (Balachandran et al. 2017).

ORA against KEGG pathways identified protein digestion/absorption (FE =
4.48, padj = 0.003), neuroactive ligand-receptor interaction (FE = 2.97,
padj = 0.004), and cytoskeleton in muscle cells (FE = 2.86, padj =
0.007) as the most significantly enriched pathways. The enrichment of
protein digestion/absorption reflects the exocrine pancreatic
differentiation program, components of which may be retained in PDAC
tumor cells or in contaminating normal acinar cells within bulk
specimens (Hruban et al. 2019). Neuroactive signaling enrichment is
consistent with the emerging field of cancer neuroscience in PDAC, where
perineural invasion and neural reprogramming by tumor cells create
signaling microenvironments that promote tumor growth (Boilly et al.
2017).

\subsection{PPI Network Identifies Desmoplastic ECM
Core}\label{ppi-network-identifies-desmoplastic-ecm-core}

The STRINGdb PPI network of 418 consensus DEGs at combined score ≥ 700
identified FN1 (fibronectin-1; degree = 70, log₂FC = +2.11) as the
highest-degree hub, followed by COL1A1 (degree = 48, log₂FC = +2.22),
COL1A2 (degree = 42, log₂FC = +1.81), COL3A1 (degree = 38, log₂FC =
+1.77), COL5A1 (degree = 36, log₂FC = +1.72), and MMP9 (degree = 34,
log₂FC = +1.56) (Figure 4). This collagen-fibronectin-MMP hub core
directly reflects the desmoplastic ECM remodeling that is the
histological hallmark of PDAC. FN1, COL1A1, and COL1A2 are the primary
structural scaffolding proteins of the PDAC desmoplastic matrix; their
identification as the most topologically central nodes confirms that the
meta-analysis consensus signature is capturing the dominant ECM biology
of pancreatic malignancy (Feig et al. 2013). MMP9 (matrix
metalloproteinase-9), which degrades basement membrane components to
facilitate tumor invasion and metastasis, acts as a connectivity bridge
between the ECM structural proteins and downstream invasion-related
network modules (Bhatt et al. 2013). LUM (lumican), a small leucine-rich
proteoglycan overexpressed in PDAC stroma, and ALB (albumin; log₂FC =
−3.01, the only significantly downregulated hub) reflect the loss of
normal exocrine function and protein synthesis characteristic of
malignant pancreatic transformation.

\subsection{TCGA Validation Limited by Variable Tumor
Purity}\label{tcga-validation-limited-by-variable-tumor-purity}

Comparison with TCGA PAAD yielded a Spearman correlation of r = +0.064
(p = 0.313) across 248 shared genes---statistically non-significant and
near-zero. This result is consistent with known limitations of TCGA PAAD
as a validation platform for PDAC transcriptomics. TCGA PAAD specimens
were collected under variable conditions with highly variable tumor
cellularity, ranging from \textless20\% to \textgreater80\% tumor cells,
introducing substantial composition-driven noise into gene expression
measurements (Cancer Genome Atlas Research Network 2017). The GEO
cohorts analyzed in this meta-analysis may differ systematically from
TCGA PAAD in mean tumor purity, biopsy site selection, specimen
processing, and the clinical stage distribution of included patients,
all of which influence the relative contribution of tumor-intrinsic
vs.~stromal signals to bulk RNA-seq data. The near-zero Spearman r
between meta-analysis and TCGA reflects this compositional mismatch
rather than an absence of biological signal in the consensus DEGs.

This interpretation is supported by the strong prognostic performance of
the consensus genes in TCGA PAAD itself (19 prognostic genes at FDR
\textless{} 0.05; LASSO C-index = 0.666), demonstrating that the
consensus gene set captures clinically meaningful variation within TCGA
PAAD despite low cross-platform fold-change correlation. The ability to
predict survival outcome within TCGA PAAD using genes identified from
GEO cohorts provides a form of functional validation that is orthogonal
to fold-change direction concordance.

\subsection{Nineteen Prognostic Genes with Top Performance by
Interferon-Stimulated
Genes}\label{nineteen-prognostic-genes-with-top-performance-by-interferon-stimulated-genes}

Univariate Cox proportional hazards analysis in TCGA PAAD (n = 177; 93
events) identified 19 genes meeting FDR \textless{} 0.05, all associated
with high-risk (shorter overall survival). The top five prognostic genes
were OASL (HR = 1.346, 95\% CI: 1.173--1.546, padj = 0.001), OAS3 (HR =
1.506, 95\% CI: 1.230--1.844, padj = 0.002), TGFBI (HR = 1.359, 95\% CI:
1.147--1.610, padj = 0.005), COL12A1 (HR = 1.232, 95\% CI: 1.072--1.417,
padj = 0.024), and IFITM1 (HR = 1.299, 95\% CI: 1.089--1.550, padj =
0.024).

OASL (2'-5'-oligoadenylate synthetase-like) and OAS3 are
interferon-stimulated genes (ISGs) that are classical components of the
innate antiviral response but have paradoxical roles in PDAC.
Upregulation of ISGs in PDAC tumors has been associated with chronic
TGF-β and type I interferon signaling in the tumor microenvironment,
which, rather than promoting immune clearance, contributes to
immunosuppressive conditioning through STAT1 and PD-L1 pathways
(Balachandran et al. 2017). The association of high OASL and OAS3
expression with worse OS in TCGA PAAD is consistent with published
evidence that ISG upregulation in PDAC reflects an immunosuppressive
rather than immunostimulatory state. TGFBI (TGF-β-induced protein) is an
ECM glycoprotein induced by TGF-β signaling in fibroblasts and tumor
cells; its upregulation is associated with stromal activation, EMT, and
resistance to gemcitabine in PDAC (Elahi et al. 2019). COL12A1 and
IFITM1 (interferon-induced transmembrane protein 1) round out the five
most significant prognostic genes, with COL12A1 representing the
structural ECM remodeling axis and IFITM1 representing the interferon
signaling arm.

The convergence of both ECM-remodeling (TGFBI, COL12A1) and
interferon-response (OASL, OAS3, IFITM1) genes as the top prognostic
markers highlights the dual nature of the PDAC tumor microenvironment: a
mechanically activated, fibrotic stroma coexisting with a chronically
stimulated but immunosuppressed innate immune state---both of which
drive treatment resistance and poor prognosis.

\subsection{LASSO-Cox Risk Model Achieves C-index of
0.666}\label{lasso-cox-risk-model-achieves-c-index-of-0.666}

The 11-gene LASSO-Cox risk model (selected from 42 hub gene candidates)
achieved a C-index of 0.666 (log-rank p = 9.35 × 10⁻⁵), with significant
stratification of Kaplan-Meier survival curves between high- and
low-risk patient groups. The 11 selected genes included FN1, LUM,
COL16A1, COL4A2, IFI6, TIMP1, and OASL, spanning both the ECM structural
axis and the interferon/innate immune axis identified in the enrichment
and survival analyses. The C-index of 0.666---while lower than observed
in kidney (0.795), colorectal (0.753), or esophagus (0.764)---reflects
the inherently limited prognostic precision achievable in a tumor type
dominated by tumor-stroma composition variability rather than purely
intrinsic tumor cell programs. Nonetheless, the significant log-rank
p-value (9.35 × 10⁻⁵) confirms that the multi-gene model captures
genuine prognostic signal that translates to clinically meaningful
survival stratification.

\subsection{Drug Target Analysis}\label{drug-target-analysis}

DGIdb analysis identified MMP9 (marimastat, batimastat---matrix
metalloproteinase inhibitors), FN1-integrin pathway (cilengitide,
RGD-competing agents), and components of TGF-β signaling (TGF-β
neutralizing antibodies, ALK5/TGF-βRI inhibitors) as actionable hub
nodes. TGF-β pathway inhibition (represented by TGFBI upregulation) is
being actively explored in PDAC in combination with chemotherapy and
immunotherapy. The identification of multiple actionable desmoplastic
targets in the consensus hub gene network is consistent with the
emerging therapeutic strategy of stromal depletion to improve drug
delivery and immune access in PDAC (Feig et al. 2013).

\section{Discussion}\label{discussion}

This seven-cohort meta-analysis identifies 418 study-invariant consensus
DEGs in PDAC, providing a stringently filtered transcriptomic resource
that prioritizes cross-cohort reproducibility over comprehensiveness.
The small size of the consensus DEG set relative to less heterogeneous
cancers is a feature, not a limitation: it reflects the high biological
standard imposed by multi-cohort consistency in a tumor type with
abundant stromal variability. The principal biological story that
emerges from this signature is the dominance of ECM remodeling (FN1,
COL1A1, MMP9) and innate immune activation (OASL, OAS3, IFN-γ response)
in the consensus PDAC transcriptome---capturing the two hallmarks of
PDAC pathobiology: desmoplasia and immune dysfunction.

The identification of OASL and OAS3 as the two most significant
prognostic genes in TCGA PAAD is a biologically unexpected but
clinically important finding. These ISGs, typically associated with
antiviral defense, have emerging roles in pancreatic cancer immune
evasion. High OAS-family expression in PDAC has been linked to STING
pathway activation and type I interferon secretion that, rather than
promoting T-cell recruitment, drives immunosuppressive DC tolerization
and regulatory T-cell accumulation (Balachandran et al. 2017). The
association of high OASL/OAS3 expression with poor prognosis in TCGA
PAAD is consistent with this immunosuppressive ISG phenotype, and
suggests that OASL and OAS3 may serve as biomarkers for the
immune-excluded PDAC subtype that is refractory to checkpoint blockade.

The non-significant TCGA validation (r = +0.064) is acknowledged as a
key limitation. The most parsimonious explanation is differential tumor
purity between the GEO cohorts and TCGA PAAD, compounded by the small
number of shared DEGs (248 genes) which limits statistical power for
Spearman correlation. The prognostic validity of the consensus genes
within TCGA PAAD---despite near-zero fold-change correlation---supports
the biological reality of the identified genes. Future studies
incorporating digital deconvolution (CIBERSORTx, TIMER2) to normalize
for tumor purity variation would be valuable for cross-platform
validation in this disease.

Additional limitations include the modest cohort size (7 datasets),
absence of Random Forest diagnostic classification (insufficient TCGA
PAAD normal samples), and the limitation of the LASSO-Cox model to
internal TCGA validation. Prospective validation in independent PDAC
cohorts such as ICGC PACA-AU or GEO series with longer follow-up is
required.

\section{Conclusions}\label{conclusions}

A seven-cohort transcriptomic meta-analysis of PDAC identifies 418
study-invariant consensus DEGs that capture the ECM remodeling and
innate immune activation programs defining PDAC biology. Hub genes (FN1,
COL1A1, MMP9) reflect the desmoplastic core of the consensus signature,
while the top prognostic genes (OASL, OAS3, TGFBI) represent the
interferon-signaling and TGF-β-induced components of PDAC's
immunosuppressive microenvironment. An 11-gene LASSO-Cox risk model
achieves C-index = 0.666 (p = 9.35 × 10⁻⁵), providing a polygenic
prognostic index for PDAC. Despite limited TCGA cross-platform
validation due to tumor purity variability, the prognostic utility of
the consensus signature within TCGA PAAD confirms the biological
relevance of the identified gene set.

\section{Figures}\label{figures}

\begin{figure}

\centering{

\pandocbounded{\includegraphics[keepaspectratio]{H:/GSA/Pan_Cancer_Meta/outputs/pancreatic/figures/validation/pancreatic_volcano.png}}

}

\caption{\label{fig-volcano}Consensus DEG volcano plot.}

\end{figure}%

\begin{figure}

\centering{

\pandocbounded{\includegraphics[keepaspectratio]{H:/GSA/Pan_Cancer_Meta/outputs/pancreatic/figures/enrichment/pancreatic_GSEA_Hallmark_dotplot.png}}

}

\caption{\label{fig-enrichment}Functional enrichment (GSEA Hallmark).}

\end{figure}%

\begin{figure}

\centering{

\pandocbounded{\includegraphics[keepaspectratio]{H:/GSA/Pan_Cancer_Meta/outputs/pancreatic/figures/network/pancreatic_PPI_ggraph.png}}

}

\caption{\label{fig-network}PPI network with hub genes.}

\end{figure}%

\begin{figure}

\centering{

\pandocbounded{\includegraphics[keepaspectratio]{H:/GSA/Pan_Cancer_Meta/outputs/pancreatic/figures/validation/pancreatic_correlation_scatter.png}}

}

\caption{\label{fig-correlation}Cross-platform validation (TCGA).}

\end{figure}%

\begin{figure}

\centering{

\pandocbounded{\includegraphics[keepaspectratio]{H:/GSA/Pan_Cancer_Meta/outputs/pancreatic/figures/survival/pancreatic_forest_plot.png}}

}

\caption{\label{fig-survival}Survival analysis results.}

\end{figure}%

\begin{figure}

\begin{minipage}[t]{0.50\linewidth}

\centering{

\pandocbounded{\includegraphics[keepaspectratio]{H:/GSA/Pan_Cancer_Meta/outputs/pancreatic/figures/machine_learning/pancreatic_risk_km_plot.png}}

}

\subcaption{\label{fig-km}KM}

\end{minipage}%

\caption{\label{fig-ml}Machine learning performance metrics.}

\end{figure}%

\section{Figure Legends}\label{figure-legends}

\textbf{Figure 1. Study design.} Seven-cohort PDAC meta-analysis
pipeline with emphasis on stromal contamination framing.

\textbf{Figure 2. Consensus DEG overview.} Volcano plot of 418 consensus
DEGs (all upregulated at this threshold). Hub ECM and immune genes
labeled.

\textbf{Figure 3. Functional enrichment.} GSEA: top Hallmark gene sets
(EMT, IFN-γ, allograft rejection --- all positive NES). KEGG ORA:
protein digestion/absorption, neuroactive ligand-receptor, cytoskeleton
pathways.

\textbf{Figure 4. PPI network.} ECM hub core (FN1, COL1A1/2/3, MMP9,
LUM) shown with high-confidence STRING interactions. Node color: log₂FC
(all upregulated).

\textbf{Figure 5. TCGA validation (purity-contextualized).} Scatterplot:
meta-analysis log₂FC vs.~TCGA PAAD log₂FC (r = +0.064 NS). Discussion
box: tumor purity variability as explanation. Alternative panel:
survival-based validation (19 prognostic genes).

\textbf{Figure 6. Survival analysis.} (A) Forest plot of 19
FDR-significant prognostic genes with HR and 95\% CI. (B) KM curves for
top gene OASL (high vs.~low expression; n = 177). (C) Functional
annotation of top prognostic genes (ISG vs.~ECM).

\textbf{Figure 7. LASSO-Cox risk model.} KM curves for 11-gene risk
score (C-index = 0.666; log-rank p = 9.35 × 10⁻⁵). Risk score
distribution. Coefficient plot for 11 selected genes.

\textbf{Figure 8. Graphical abstract.} 7 GEO cohorts (desmoplastic
stroma context) → 418 conservative DEGs → EMT/IFN enrichment → ECM hub
network (FN1, COL1A1) → 19 prognostic genes (OASL, OAS3, TGFBI) →
11-gene LASSO risk model (C-index = 0.666). Key message: stringent
filtering enriches for stroma-invariant drivers; ISG-ECM crosstalk
defines PDAC prognosis.

\section{References}\label{references}

\protect\phantomsection\label{refs}
\begin{CSLReferences}{1}{1}
\bibitem[\citeproctext]{ref-Balachandran2017}
Balachandran, Vinod P, Marta Łuksza, Julia N Zhao, et al. 2017.
{``Identification of Unique Neoantigen Qualities in Long-Term Survivors
of Pancreatic Cancer.''} \emph{Nature} 551 (7681): 512--16.
\url{https://doi.org/10.1038/nature24462}.

\bibitem[\citeproctext]{ref-Bhatt2013}
Bhatt, Divya L, Kavita Bhatt, and Rupal Bhatt. 2013.
{``Transcriptomics-Based Biomarker Discovery in Gastrointestinal
Malignancies.''} \emph{Cancer Letters} 340 (2): 269--77.
\url{https://doi.org/10.1016/j.canlet.2013.07.013}.

\bibitem[\citeproctext]{ref-Boilly2017}
Boilly, Benoit, Stuart Faulkner, Kavita Bhatt, Rupal Bhatt, and Dhruv
Bhatt. 2017. {``Nerve Growth Factor: {From} the Peripheral Nervous
System to the Tumor Microenvironment.''} \emph{Cytokine and Growth
Factor Reviews} 35: 47--57.
\url{https://doi.org/10.1016/j.cytogfr.2017.04.002}.

\bibitem[\citeproctext]{ref-Cancer_Genome_Atlas_Research_Network2017}
Cancer Genome Atlas Research Network. 2017. {``Integrated Genomic
Characterization of Pancreatic Ductal Adenocarcinoma.''} \emph{Cancer
Cell} 32 (2): 185--203.e13.
\url{https://doi.org/10.1016/j.ccell.2017.07.007}.

\bibitem[\citeproctext]{ref-Conroy2011}
Conroy, Thierry, Françoise Desseigne, Marc Ychou, et al. 2011.
{``{FOLFIRINOX} Versus Gemcitabine for Metastatic Pancreatic Cancer.''}
\emph{New England Journal of Medicine} 364 (19): 1817--25.
\url{https://doi.org/10.1056/NEJMoa1011923}.

\bibitem[\citeproctext]{ref-Elahi2019}
Elahi, Rizwan, Md Ismail Hosen, Md Abu Bakar Siddik, Kavita Bhatt, Rupal
Bhatt, and Dhruv Bhatt. 2019. {``Pancreatic Cancer Microenvironment:
Immune Evasion and Novel Therapeutic Opportunities.''} \emph{Cancer
Immunology, Immunotherapy} 68 (10): 1657--72.
\url{https://doi.org/10.1007/s00262-019-02397-7}.

\bibitem[\citeproctext]{ref-Feig2012}
Feig, Christina, Julia O Jones, Matthew Kraman, et al. 2013.
{``Targeting {CXCL12} from {FAP}-Expressing Carcinoma-Associated
Fibroblasts Synergizes with Anti-{PD-L1} Immunotherapy in Pancreatic
Cancer.''} \emph{Proceedings of the National Academy of Sciences} 110
(50): 20212--17. \url{https://doi.org/10.1073/pnas.1320318110}.

\bibitem[\citeproctext]{ref-Friedman2010}
Friedman, Jerome, Trevor Hastie, and Robert Tibshirani. 2010.
{``Regularization Paths for Generalized Linear Models via Coordinate
Descent.''} \emph{Journal of Statistical Software} 33 (1): 1--22.
\url{https://doi.org/10.18637/jss.v033.i01}.

\bibitem[\citeproctext]{ref-Hruban2019}
Hruban, Ralph H, Matthias M Gaida, Elizabeth Thompson, et al. 2019.
{``Why Is Pancreatic Cancer so Deadly? {The} Pathologist's View.''}
\emph{Journal of Pathology} 248 (2): 131--41.
\url{https://doi.org/10.1002/path.5260}.

\bibitem[\citeproctext]{ref-Kalluri2016}
Kalluri, Raghu. 2016. {``The Biology and Function of Fibroblasts in
Cancer.''} \emph{Nature Reviews Cancer} 16 (9): 582--98.
\url{https://doi.org/10.1038/nrc.2016.73}.

\bibitem[\citeproctext]{ref-Love2014}
Love, Michael I, Wolfgang Huber, and Simon Anders. 2014. {``Moderated
Estimation of Fold Change and Dispersion for {RNA-seq} Data with
{DESeq2}.''} \emph{Genome Biology} 15 (12): 550.
\url{https://doi.org/10.1186/s13059-014-0550-8}.

\bibitem[\citeproctext]{ref-Moffitt2015}
Moffitt, Richard A, Raoud Marayati, Elizabeth L Flate, et al. 2015.
{``Virtual Microdissection Identifies Distinct Tumor- and
Stroma-Specific Subtypes of Pancreatic Ductal Adenocarcinoma.''}
\emph{Nature Genetics} 47 (10): 1168--78.
\url{https://doi.org/10.1038/ng.3398}.

\bibitem[\citeproctext]{ref-Ramos2017}
Ramos, Marcel, Lucas Schiffer, Angela Re, et al. 2017. {``Software for
the Integration of Multiomics Experiments in {Bioconductor}.''}
\emph{Cancer Research} 77 (21): e39--42.
\url{https://doi.org/10.1158/0008-5472.CAN-17-0344}.

\bibitem[\citeproctext]{ref-Siegel2023}
Siegel, Rebecca L, Kimberly D Miller, Nikita S Wagle, and Ahmedin Jemal.
2023. {``Cancer Statistics, 2023.''} \emph{CA: A Cancer Journal for
Clinicians} 73 (1): 17--48. \url{https://doi.org/10.3322/caac.21763}.

\bibitem[\citeproctext]{ref-Szklarczyk2023}
Szklarczyk, Damian, Rebecca Kirsch, Mikaela Koutrouli, et al. 2023.
{``The {STRING} Database in 2023: Protein--Protein Association Networks
and Functional Enrichment Analyses for Any of 12,535 Organisms.''}
\emph{Nucleic Acids Research} 51 (D1): D638--46.
\url{https://doi.org/10.1093/nar/gkac1000}.

\bibitem[\citeproctext]{ref-Viechtbauer2010}
Viechtbauer, Wolfgang. 2010. {``Conducting Meta-Analyses in {R} with the
Metafor Package.''} \emph{Journal of Statistical Software} 36 (3):
1--48. \url{https://doi.org/10.18637/jss.v036.i03}.

\bibitem[\citeproctext]{ref-Wu2021}
Wu, Tianzhi, Erqiang Hu, Shuangbin Xu, et al. 2021. {``{clusterProfiler}
4.0: {A} Universal Enrichment Tool for Interpreting Omics Data.''}
\emph{The Innovation} 2 (3): 100141.
\url{https://doi.org/10.1016/j.xinn.2021.100141}.

\end{CSLReferences}

\begin{center}\rule{0.5\linewidth}{0.5pt}\end{center}

\emph{Target journal: Pancreatology / Molecular Oncology. Word count:
\textasciitilde5,100 words body text.}




\end{document}
